\documentclass[tikz,border=6pt]{standalone}
\usepackage[no-math]{fontspec}
\usepackage[brazil]{babel}
\usepackage{amsmath,amssymb}
\IfFileExists{C:/Windows/Fonts/Lato_Regular.ttf}{%
  \setmainfont{Lato}[
    Path = C:/Windows/Fonts/,
    UprightFont = Lato_Regular.ttf,
    BoldFont = Lato_Bold.ttf,
    ItalicFont = Lato_Italic.ttf,
    BoldItalicFont = Lato_BoldItalic.ttf
  ]
  \setsansfont{Lato}[
    Path = C:/Windows/Fonts/,
    UprightFont = Lato_Regular.ttf,
    BoldFont = Lato_Bold.ttf,
    ItalicFont = Lato_Italic.ttf,
    BoldItalicFont = Lato_BoldItalic.ttf
  ]
}{%
  \RequirePackage[default]{lato}
}
\renewcommand{\familydefault}{\sfdefault}

\usepackage{xcolor}
\usepackage{tikz}
\usepackage{pgfplots}
\pgfplotsset{compat=1.18}
\usetikzlibrary{arrows.meta,positioning,shapes.geometric,fit,backgrounds,calc,decorations.pathreplacing,patterns}

% Paleta institucional UEFS / CHF662
\definecolor{uefsNavy}{HTML}{1B2A4A}
\definecolor{uefsTeal}{HTML}{0E7C7B}
\definecolor{uefsForest}{HTML}{2D6A4F}
\definecolor{uefsGold}{HTML}{C4983B}
\definecolor{uefsGoldDark}{HTML}{8F6B1E}
\definecolor{uefsAlert}{HTML}{B02A37}
\definecolor{uefsOrange}{HTML}{D95F02}
\definecolor{uefsInk}{HTML}{22252A}
\definecolor{uefsMuted}{HTML}{5C6470}
\definecolor{uefsWater}{HTML}{2E86AB}
\definecolor{uefsSand}{HTML}{D9B382}
\definecolor{bgNavy}{HTML}{F2F5F9}
\definecolor{bgTeal}{HTML}{E9F5F4}
\definecolor{bgForest}{HTML}{EEF5F0}
\definecolor{bgGold}{HTML}{FCF7EB}
\definecolor{bgAlert}{HTML}{FDF2F2}
\definecolor{bgGray}{HTML}{F8FAFC}
\definecolor{lineGray}{HTML}{CFD8DC}

\tikzset{
  >=Stealth,
  aaCard/.style={draw=#1, fill=white, rounded corners=4pt, line width=1pt, inner sep=6pt, align=center, font=\small},
  aaBox/.style={draw=#1, rounded corners=4pt, line width=1.1pt, inner sep=6pt, align=center, font=\footnotesize\sffamily, text=uefsInk},
  aaTitle/.style={font=\normalsize\bfseries\sffamily, text=uefsNavy, align=center},
  aaSub/.style={font=\footnotesize\sffamily, text=uefsMuted, align=center},
  aaArrow/.style={->, >=Stealth, line width=1.4pt},
  aaDArrow/.style={<->, >=Stealth, line width=1.4pt}
}


\begin{document}
\begin{tikzpicture}

\def\cardW{3.85}
\def\cardH{3.45}
\def\hdrH{0.92}
\def\gap{0.75}

% 1. Ação do Projeto
\pgfmathsetmacro{\xa}{0}
\begin{scope}
  \clip[rounded corners=4pt] (\xa,0) rectangle (\xa+\cardW,\cardH);
  \fill[bgNavy] (\xa,0) rectangle (\xa+\cardW,\cardH);
  \fill[uefsNavy] (\xa,\cardH-\hdrH) rectangle (\xa+\cardW,\cardH);
\end{scope}
\draw[uefsNavy, line width=1.1pt, rounded corners=4pt] (\xa,0) rectangle (\xa+\cardW,\cardH);
\node[text width=3.45cm, align=center, font=\footnotesize\bfseries\sffamily, text=white]
  at (\xa+\cardW/2, \cardH-\hdrH/2) {Ação do Projeto};
\node[text width=3.35cm, align=center, font=\footnotesize\sffamily, text=uefsInk]
  at (\xa+\cardW/2, {(\cardH-\hdrH)/2})
  {Abertura de estrada em encosta sem drenagem pluvial};

% 2. Impacto Direto (1ª ordem)
\pgfmathsetmacro{\xb}{\cardW+\gap}
\begin{scope}
  \clip[rounded corners=4pt] (\xb,0) rectangle (\xb+\cardW,\cardH);
  \fill[bgNavy] (\xb,0) rectangle (\xb+\cardW,\cardH);
  \fill[uefsWater] (\xb,\cardH-\hdrH) rectangle (\xb+\cardW,\cardH);
\end{scope}
\draw[uefsWater, line width=1.1pt, rounded corners=4pt] (\xb,0) rectangle (\xb+\cardW,\cardH);
\node[text width=3.45cm, align=center, font=\footnotesize\bfseries\sffamily, text=white]
  at (\xb+\cardW/2, \cardH-\hdrH/2) {Impacto Direto (1\textsuperscript{a} ordem)};
\node[text width=3.35cm, align=center, font=\footnotesize\sffamily, text=uefsInk]
  at (\xb+\cardW/2, {(\cardH-\hdrH)/2})
  {Remoção da Caatinga e compactação ou exposição do horizonte superficial};

% 3. Impacto Indireto (2ª ordem)
\pgfmathsetmacro{\xc}{2*(\cardW+\gap)}
\begin{scope}
  \clip[rounded corners=4pt] (\xc,0) rectangle (\xc+\cardW,\cardH);
  \fill[bgTeal] (\xc,0) rectangle (\xc+\cardW,\cardH);
  \fill[uefsTeal] (\xc,\cardH-\hdrH) rectangle (\xc+\cardW,\cardH);
\end{scope}
\draw[uefsTeal, line width=1.1pt, rounded corners=4pt] (\xc,0) rectangle (\xc+\cardW,\cardH);
\node[text width=3.45cm, align=center, font=\footnotesize\bfseries\sffamily, text=white]
  at (\xc+\cardW/2, \cardH-\hdrH/2) {Impacto Indireto (2\textsuperscript{a} ordem)};
\node[text width=3.35cm, align=center, font=\footnotesize\sffamily, text=uefsInk]
  at (\xc+\cardW/2, {(\cardH-\hdrH)/2})
  {Escoamento concentrado, sulcos, ravinas e carreamento de sedimentos};

% 4. Impacto Indireto (3ª ordem)
\pgfmathsetmacro{\xd}{3*(\cardW+\gap)}
\begin{scope}
  \clip[rounded corners=4pt] (\xd,0) rectangle (\xd+\cardW,\cardH);
  \fill[bgGold] (\xd,0) rectangle (\xd+\cardW,\cardH);
  \fill[uefsOrange] (\xd,\cardH-\hdrH) rectangle (\xd+\cardW,\cardH);
\end{scope}
\draw[uefsOrange, line width=1.1pt, rounded corners=4pt] (\xd,0) rectangle (\xd+\cardW,\cardH);
\node[text width=3.45cm, align=center, font=\footnotesize\bfseries\sffamily, text=white]
  at (\xd+\cardW/2, \cardH-\hdrH/2) {Impacto Indireto (3\textsuperscript{a} ordem)};
\node[text width=3.35cm, align=center, font=\footnotesize\sffamily, text=uefsInk]
  at (\xd+\cardW/2, {(\cardH-\hdrH)/2})
  {Assoreamento de açude comunitário e escassez hídrica na estiagem};

% Setas entre cartões
\draw[aaArrow, color=uefsNavy] (\xa+\cardW+0.08, \cardH/2) -- (\xb-0.08, \cardH/2);
\draw[aaArrow, color=uefsWater] (\xb+\cardW+0.08, \cardH/2) -- (\xc-0.08, \cardH/2);
\draw[aaArrow, color=uefsTeal] (\xc+\cardW+0.08, \cardH/2) -- (\xd-0.08, \cardH/2);

\end{tikzpicture}
\end{document}
